\scanpage{57}
% manuscript line 1
De plus, si $u,v\in L^{(2)}(\deleted{\unsure{F}}E,\deleted{\unsure{E}}F)$, on a

% manuscript line 2
\[
\boxed{(u,v)=\operatorname{Tr}uv^*=\operatorname{Tr}v^*u} \tag*{\textcircled{25}}
\]

% manuscript line 3
\uline{Démonstration} Pour prouver (\unsure{25|24}), il suffit \deleted{\unsure{pour raison de continuité|par raison de continuité}}
\deleted{\unsure{si}} de se borner au cas où $u,v$ sont de rang fini
(alors \deleted{l'application} par prolongement par continuité on aura
une application bilinéaire de norme $\leq1$ de $L^2(E,F)\times
L^2(F,G)$ dans $\unsure{L^2(E,G)|L^1(E,G)}$, qui n'est \unsure{par} \deleted{\unsure{donc}} ne pourra
être que $(u,v)\to vu$). Or, \deleted{\unsure{ou|on}} d'après (22),
on est ramené \deleted{\unsure{al}} à prouver $|\operatorname{Tr} Avu|\leq\|v\|_2\|u\|_2$
pour \unsure{$\forall$} $A\in L(G,E)$, $\|A\|\leq1$. Comme $\|Av\|_2\leq\|A\|\|v\|_2$
\deleted{(prop. \unsure{1})} (formule (17)) on est ramené à prouver
$|\operatorname{Tr}wu|\leq\|w\|_2\|u\|_2$. En vertu de $\|v\|=\|v^*\|$, on
peut \deleted{\unsure{remplacer $v$ par $v^*$} (\illegible{1})} \inserted{poser $w=v^*$ \deleted{\unsure{$(v\in\bar E\otimes F)$}} $(v\in E'\otimes F)$}, et on est ramené
à prouver $|\operatorname{Tr}v^*u|\leq\|v^*\|_2\|u\|_2$. Mais cela résultera
de la formule (25), puisque $|(u,v)|\leq\|u\|_2\|v\|_2$.
\deleted{\unsure{Par|Ou}} Or la formule (25) pour opérateurs de rang
fini est triviale \,: il suffit de la vérifier pour
$u=\bar a_{\deleted{\unsure{1}}}\otimes b_{\deleted{\unsure{1}}}$, $v=\overline{a'}\otimes b'$, d'où $v^*=\overline{b'}\otimes a'$,
$v^*u=(b,\unsure{\overline{b'}|b'})\,\bar a\otimes a'$, \deleted{$\unsure{uv^*=(a,a')\bar b'\otimes b|uv^*=(a',a)\bar b\otimes b}$} d'où
\deleted{\unsure{Tr}} $(u,v)=(\bar a,\overline{a'})(b,b')=\deleted{\unsure{(a,a')|(a^*,a')}}\inserted{(a',a)}(b,b')$, $\operatorname{Tr}v^*u=(a',a)(b,b')$.
Comme $\operatorname{Tr}v^*u=\operatorname{Tr}uv^*$ (\inserted{formule de} commutation pour la trace)
(25) est démontrée pour $u,v$ de rang fini, \deleted{\unsure{Par là}} Ainsi
la première partie du th. est démontré, donc aussi
la formule (25) pour \deleted{tout} $u\in L^2(E,F)$, $v\in L^2(E,F)$, par
prolongement par continuité.

% manuscript line 26
% archive sheet number: 58
